% -----------------------------------------------------------------------------------------------------
\section{超领头对数、相干性破坏与“因子化恢复”}\label{sec:front-sll}

在强子对撞的非全局观测量（喷注之间的间隙、喷注否决）中，否决标度 $Q_0\ll Q$，$L=\ln(Q/Q_0)$，期望的领头对数是 $\alpha_s^nL^n$。硬振幅中的复 Coulomb（Glauber）相位破坏色相干性，产生\emph{双}对数 $\alpha_s^3L^3(\alpha_sL^2)^n$，在四圈即 $\alpha_s^4L^5$：这就是超领头对数（super-leading logarithms，SLL）。由于 PDF 演化是单对数的，若 PDF 因子化要成立，这种双对数的共线演化必须被入射部分子之间的低能相互作用补偿。

\begin{proposition}[非全局 LHC 观测量的因子化与超领头对数的全阶重求和 \tagA]\label{prop:front-sll-all-orders}
对否决标度为 $Q_0\ll Q$ 的强子对撞非全局观测量，截面因子化为
\begin{equation}\label{eq:front-hw}
\sigma(Q_0)=\sum_{a_1a_2}\int\dd x_1\dd x_2\sum_m\bigl\langle\mathcal H_m(\{n\},Q,\mu)\otimes\mathcal W_m(\{n\},Q_0,x_1,x_2,\mu)\bigr\rangle ,
\end{equation}
即硬函数与否决标度上的低能矩阵元之积（在多重数空间与色空间中的矩阵，$\otimes$ 为对方向的积分，$\langle\cdot\rangle$ 为色迹）；单圈反常维数含软与共线两部分，共线部分带正比于 $\ii\pi$ 的 Glauber 相位（\cref{prop:sudakov-coulomb}），二者不对易产生 SLL，并可在领头对数精度上对任意 $2\to M$ 过程全阶重求和：所得级数交错，各项远大于其和；对 $pp\to2$ 喷注显著，对 $H/Z(+$喷注$)$ 受压低（\cite{Becher:2023mtx}，据该文摘要）。含任意多个相位 $(\ii\pi)^2$ 的完整“Glauber 级数”在领头对数精度上见\cite{Boer:2023jsy,Boer:2023ljq,Boer:2024hzh,Boer:2024xzy}（后两篇分别含跑动耦合与大 $N_c$），$pp\to2$ 喷注的全部部分子道见\cite{Becher:2024nqc}（某些 $q\bar q\to q\bar q$ 干涉项只受 $1/N_c$ 压低），有质量末态带来额外的 Coulomb 相位与新的 SLL 来源\cite{Banerjee:2025kkq}（均据摘要）。SLL 由 Forshaw、Kyrieleis 与 Seymour 发现\cite{ForshawKyrieleisSeymour2006}，色基无关性见\cite{Forshaw:2008cq}，排序变量见\cite{Angeles-Martinez:2015rna}，几乎一切软敏感强子对撞观测量中的相干性破坏见\cite{Forshaw:2021fxs}（含 2024 年勘误）。（本命题只涉及部分子层次的低能矩阵元与硬函数。）
% errtype: E-S
\uses{thm:jets-ngl-factorization, prop:sudakov-coulomb}
\end{proposition}

\begin{proofsketch}
用 SCET 推导\cref{eq:front-hw}（其 $e^+e^-$ 版本是\cref{thm:jets-ngl-factorization}），算出 $\mathcal H_m$、$\mathcal W_m$ 的单圈反常维数矩阵，并把重整化群演化算符按 Glauber 相位插入的次数展开、在领头对数精度上逐项求和。完整论证见\cite{Becher:2023mtx}。
\end{proofsketch}

\paragraph{地位与缺口。}[A]：对 $\mathcal H\otimes\mathcal W$ 结构与领头对数重求和（有效理论推导，单圈重整化群）。否决标度以下的 PDF 因子化是另一个问题（\cref{prop:front-restoration}，[P]）。核心陈述是\cref{prop:jets-sll}（[P]，SLL 的存在）。

\begin{proposition}[超领头对数的塔与 $q\bar q$ 单态交换的闭式 \tagA]\label{prop:front-sll-closed-form}
迭代\cref{prop:front-sll-all-orders}的单圈反常维数，得到 SLL 的塔
\begin{equation}\label{eq:front-sll-tower}
\hat\sigma_n^{\rm SLL}=\Bigl(\frac{\alpha_s}{4\pi}\Bigr)^{n+3}L^{2n+3}\,\frac{(-4)^n\,n!}{(2n+3)!}\sum_{r=0}^n\frac{(2r)!}{4^r(r!)^2}\,C_{rn},
\end{equation}
色迹 $C_{rn}$ 显式已知（\cite[式~(14)--(21)]{Becher:2021zkk}，作者用 ColorMath 核对到八圈）。对夸克—夸克散射中 $t$ 道 $3\otimes\bar3$ 色单态的交换（如光子交换），在前向散射中（$J_2=J_{43}/2=\Delta Y$，$\Delta Y$ 为间隙的快度宽度），
\begin{equation}\label{eq:front-sll-closed}
\Delta\hat\sigma^{(S)}=\sum_{n\ge0}\hat\sigma_n^{\rm SLL}=-\hat\sigma_B\,\frac{4C_F}{3\pi}\,\alpha_s^3L^3\,\Delta Y\;{}_2F_2\Bigl(1,1;2,\frac52;-w\Bigr),\qquad w=\frac{N_c\alpha_s}{\pi}L^2 ,
\end{equation}
$\hat\sigma_B$ 为 Born 截面；$n=0$ 项是单对数的，第一个 SLL 为 $\hat\sigma_1^{(S)}=\hat\sigma_B\frac{4}{15\pi^2}C_FN_c\alpha_s^4L^5\Delta Y$；大 $w$ 时 $w\,{}_2F_2(1,1;2,\frac52;-w)/\ln w\to\frac32$（$w=10^8$ 时为 $1.497$），是缓慢的幂律式衰减而非 Sudakov 型的 $\ee^{-cw}$。SLL 也出现在 $2\to1$ 过程（从四圈起）与 $2\to0$ 过程（如带中心喷注否决的 $q\bar q\to V$，从五圈起）中。（本命题只涉及部分子截面。）
% VF-FRO-7 (the closed form is for the colour-singlet exchange channel of quark-quark scattering, Becher:2021zkk eq. (19), (21))
% errtype: E-S
\uses{prop:front-sll-all-orders}
\end{proposition}

\begin{proofsketch}
塔\cref{eq:front-sll-tower}来自迭代的反常维数（引自\cite{Becher:2021zkk}）；闭式是级数恒等式，下面完整给出。单态道的色迹为\cite[式~(19)]{Becher:2021zkk}
\begin{equation*}
C^{(S)}_{rn}=\hat\sigma_B\,2^{8-r}\pi^2(4N_c)^nC_F\bigl[-J_{43}+2J_2(1-\delta_{r0})\bigr];
\end{equation*}
前向散射中 $r\ge1$ 的项为零，只剩 $C^{(S)}_{0n}=-\hat\sigma_B2^9\pi^2(4N_c)^nC_F\Delta Y$，代入\cref{eq:front-sll-tower}得
\begin{equation*}
\hat\sigma^{(S)}_n=-\hat\sigma_B\,\frac{8C_F}{\pi}\,\alpha_s^3L^3\Delta Y\,(-4w)^n\frac{n!}{(2n+3)!}.
\end{equation*}由倍元公式 $(2n+3)!=3\cdot2^{2n+1}(n+1)!\,(\frac52)_n$，$(-4w)^n\frac{n!}{(2n+3)!}=\frac16\frac{(-w)^n}{(n+1)(\frac52)_n}$，而 $\frac{(1)_n(1)_n}{(2)_n(\frac52)_n\,n!}=\frac{1}{(n+1)(\frac52)_n}$ 正是 ${}_2F_2(1,1;2,\frac52;-w)$ 的展开系数，求和即得\cref{eq:front-sll-closed}，对一切 $n$ 成立。大 $w$ 行为来自 ${}_2F_2(1,1;2,\frac52;-w)=\int_0^1\dd t\,{}_1F_1(1;\frac52;-wt)$ 与 ${}_1F_1(1;\frac52;-z)\simeq\frac{3}{2z}$。
% check: writing/work/W27/s_series.py (A1)-(A11); frontier check OTHERS s4 (n <= 8)
\end{proofsketch}

\paragraph{地位与缺口。}[A]：继承\cref{prop:front-sll-all-orders}；上面的级数恒等式是完整的。\cite{Becher:2021zkk}指出，取 $n=1$ 即重现 Forshaw、Kyrieleis 与 Seymour 的结果\cite{ForshawKyrieleisSeymour2006}。

\begin{example}[由闭式读出第一个超领头对数]\label{ex:front-sll-first}
${}_2F_2(1,1;2,\frac52;-w)=\sum_k\frac{(-w)^k}{(k+1)(\frac52)_k}=1-\frac w5+\frac{4w^2}{105}-\cdots$。$k=0$ 项给出单对数的 $-\hat\sigma_B\frac{4C_F}{3\pi}\alpha_s^3L^3\Delta Y$；$k=1$ 项给出 $-\hat\sigma_B\frac{4C_F}{3\pi}\alpha_s^3L^3\Delta Y\cdot(-\frac15\frac{N_c\alpha_s}{\pi}L^2)=\hat\sigma_B\frac{4}{15\pi^2}C_FN_c\alpha_s^4L^5\Delta Y$，即四圈的第一个 SLL：每多一圈多出两个对数。
\uses{prop:front-sll-closed-form}
\end{example}

\begin{proposition}[Glauber 胶子引起的“因子化恢复” \tagP]\label{prop:front-restoration}
在三圈、领头对数精度上，间隙截面的低能矩阵元 $\mathcal W_m$ 含有一个微扰的主动—主动 Glauber 贡献，它在纯维数正规化中有明确定义，并在领头对数精度上恰好具有把双对数演化变回否决标度以下单对数 DGLAP 演化（\cref{thm:pdf-dglap}）所需的结构；它“breaks soft-collinear factorization at the cross-section level but restores PDF factorization”\cite{Becher:2024kmk}。由有效理论重整化条件导出的相容关系已为三圈显式计算验证（\cite{Becher:2025igg}，据该文摘要）。（本命题只涉及部分子层次的低能矩阵元。）
\uses{prop:front-sll-all-orders, thm:pdf-dglap}
\end{proposition}

\paragraph{地位与缺口。}[P]（三圈，领头对数结构）。作者写道：“our analysis does not amount to a proof of PDF factorization, [but] demonstrates that the breaking of collinear factorization does not necessarily translate into a breaking of PDF factorization”。所涉 Glauber 区域会被自动的区域搜索程序漏掉（\cref{prop:front-hidden-regions}）。“间隙截面的因子化已被证明破坏”是过度表述：振幅层次的破坏是 [V]，间隙截面的 PDF 因子化是带恢复证据的猜想。

\begin{proposition}[1-jettiness 中的相干性破坏对数 \tagV]\label{prop:front-jettiness-cvl}
对色单态加一个喷注的产生（研究了 $qg\to qZ$）中的 1-jettiness $\tau_1=\frac1Q\sum_k\min\{n_a\cdot p_k,n_b\cdot p_k,n_J\cdot p_k\}$\cite{Stewart:2010tn}，相干性破坏对数相对于 Born 截面从 $\alpha_s^4\ln^6(1/\tau_1)$ 开始：
\begin{equation}\label{eq:front-cvl}
\begin{gathered}
\frac{\dd\sigma_1^{\rm CVL}}{\dd x_a\,\dd x_b\,\dd\mathcal B}\simeq\sum_{a=q,g}\mathcal A_a\Bigl(\frac{\alpha_s}{\pi}\Bigr)^4(-\ii\pi)^2\frac{1}{480}\ln^6\frac1{\tau_1},\\
\mathcal A_q=\frac83N_c^2T_R^4f^q_Af^g_B\sigma^{(0)}_{qg},\qquad \mathcal A_g=-\frac{32}3N_c^2T_R^4f^g_Af^q_B\sigma^{(0)}_{gq}
\end{gathered}
\end{equation}
（对 $\tau_1$ 的累积分布、对 Born 变量微分；PDF 取在 $\mu_F=\tau_1Q$；固定耦合），并预期一般为 $\alpha_s^nL^{2n-2}$（$n\ge4$）。因此累积分布不能写成标准的指数化形式 $\sigma(\tau_1)=\sigma\,\ee^{-\frac{\alpha_s}{2\pi}L^2(2C_F+C_A)}[g_0(\alpha_sL)+\alpha_sg_1(\alpha_sL)+\cdots]$\cite[式~(1)--(18)]{Banfi:2025mra}。机制：Coulomb 相位使 $[\mathbf T^2_{ad},\mathbf T^2_s]\ne0$（入射部分子 $a$ 的色与其共线辐射的色被分辨开，初态角度排序被破坏）；多出的对数来自从不接触共线极限的软宽角双对数相空间面积，因而不产生额外的共线极点，与 $\tau_1Q$ 处的 PDF 因子化相容。
% VF-FRO-8
\uses{def:jets-algorithms, prop:sudakov-coulomb}
\end{proposition}

\paragraph{地位与缺口。}对强子对撞中 $\tau_1$ 的标准因子化（三条 Wilson 线的软函数）与指数化的全阶有效性为 [V]，证据：显式的领头对数相干性破坏计算（近似：强有序的软—共线发射、测量函数的多重不相关发射近似、偶极 $k_T$ 排序——$\alpha_s^4L^6$ 项与排序无关，$\alpha_s^4L^5$ 项则否）。这与\cite{Alioli:2023rxx}关于 $V+$喷注中 $\tau_1$“does not suffer from superleading or nonglobal logarithms”的说法（N$^3$LL 重求和与 $N$-jettiness 切片的基础）在超领头部分相矛盾，非全局对数不在争议之内。这一结果由两个研究组用不同方法建立：Becher 等人在含 Glauber 与新动量模式的 SCET 中重新得到了 $qg\to q$ 道的这一结果\cite{Becher:2026kbr}，并指出现有的多数全局 LHC 观测量的因子化公式需要修正。
